\documentclass[11pt]{article}

\usepackage[T1]{fontenc}
\usepackage{lmodern}
\usepackage{amsmath,amssymb,amsthm,mathtools}
\usepackage{booktabs}
\usepackage{array}
\usepackage{geometry}
\usepackage{microtype}
\usepackage[hidelinks]{hyperref}
\usepackage{aliascnt}
\usepackage[nameinlink,noabbrev]{cleveref}

\geometry{margin=1.08in}
\setlength{\parindent}{0pt}
\setlength{\parskip}{0.55em}
\allowdisplaybreaks

\newtheorem{theorem}{Theorem}[section]
\newaliascnt{proposition}{theorem}
\newtheorem{proposition}[proposition]{Proposition}
\aliascntresetthe{proposition}
\newaliascnt{lemma}{theorem}
\newtheorem{lemma}[lemma]{Lemma}
\aliascntresetthe{lemma}
\newaliascnt{corollary}{theorem}
\newtheorem{corollary}[corollary]{Corollary}
\aliascntresetthe{corollary}
\theoremstyle{remark}
\newaliascnt{remark}{theorem}
\newtheorem{remark}[remark]{Remark}
\aliascntresetthe{remark}

\crefname{theorem}{theorem}{theorems}
\crefname{proposition}{proposition}{propositions}
\crefname{lemma}{lemma}{lemmas}
\crefname{corollary}{corollary}{corollaries}
\crefname{remark}{remark}{remarks}

\newcommand{\Z}{\mathbb Z}
\newcommand{\R}{\mathbb R}
\newcommand{\Norm}{\operatorname{N}}

\title{Infinitely Many Integer Solutions of\\[0.25em]
\texorpdfstring{$x^3+x^2y^2+z^2+1=0$}{x^3 + x^2 y^2 + z^2 + 1 = 0}}
\author{}
\date{}
\hypersetup{
  pdftitle={Infinitely Many Integer Solutions of x3 + x2 y2 + z2 + 1 = 0},
  pdfauthor={}
}

\begin{document}
\maketitle
\vspace{-2.5em}

\begin{abstract}
We prove, unconditionally and by an entirely elementary argument, that
\[
 x^3+x^2y^2+z^2+1=0
\]
has infinitely many integer solutions.  The construction combines a norm identity in the Gaussian integers with an explicitly verified polynomial identity.  A self-contained, congruence-controlled Pell argument then shows that an associated quadratic polynomial assumes square values at infinitely many integer arguments.  Every numerical identity used in the construction is checked coefficient by coefficient.
\end{abstract}

\section{Statement and preliminary reduction}

We establish the following result.

\begin{theorem}[Main theorem]\label{thm:main}
The Diophantine equation
\begin{equation}\label{eq:original}
 x^3+x^2y^2+z^2+1=0
\end{equation}
has infinitely many pairwise distinct solutions $(x,y,z)\in\Z^3$.
\end{theorem}

A solution of \eqref{eq:original} necessarily has $x<0$, since all four terms on the left are nonnegative when $x\ge 0$.  Writing
\[
 a=-x>0,
\]
we obtain the equivalent equation
\begin{equation}\label{eq:reduced}
 z^2+1=a^2(a-y^2).
\end{equation}
In particular, every solution satisfies $a>y^2$.  Our task is therefore to construct infinitely many triples $(a,y,z)$ satisfying \eqref{eq:reduced}.

The signs of $y$ and $z$ are immaterial in \eqref{eq:original}.  Once one solution is found, the four sign choices $(x,\pm y,\pm z)$ are again solutions.  We shall nevertheless prove infinitude already at the level of the $x$-coordinates.

\section{A Gaussian norm construction}

The basic algebraic mechanism is the following lemma.

\begin{lemma}[Gaussian norm mechanism]\label{lem:gaussian}
Let $P,Q,R,S,Y\in\Z$ satisfy
\begin{align}
 2PQR+(P^2-Q^2)S&=1, \label{eq:bezout-condition}\\
 P^2+Q^2-R^2-S^2&=Y^2. \label{eq:square-difference}
\end{align}
Define
\begin{align}
 X&=-(P^2+Q^2), \label{eq:X-definition}\\
 Z&=(P^2-Q^2)R-2PQS. \label{eq:Z-definition}
\end{align}
Then
\[
 X^3+X^2Y^2+Z^2+1=0.
\]
\end{lemma}

\begin{proof}
Put
\[
 A_0=P^2-Q^2,
 \qquad
 B_0=2PQ.
\]
Condition \eqref{eq:bezout-condition} says that
\[
 B_0R+A_0S=1,
\]
whereas \eqref{eq:Z-definition} says that
\[
 Z=A_0R-B_0S.
\]
The elementary two-square identity gives
\begin{align*}
 Z^2+1
 &= (A_0R-B_0S)^2+(B_0R+A_0S)^2\\
 &= (A_0^2+B_0^2)(R^2+S^2).
\end{align*}
Now
\[
 A_0^2+B_0^2
 =(P^2-Q^2)^2+4P^2Q^2
 =(P^2+Q^2)^2.
\]
Writing $a=P^2+Q^2$, condition \eqref{eq:square-difference} becomes
\[
 R^2+S^2=a-Y^2.
\]
Consequently,
\begin{equation}\label{eq:norm-conclusion}
 Z^2+1=a^2(a-Y^2).
\end{equation}
Condition \eqref{eq:bezout-condition} also shows that $P$ and $Q$ cannot both vanish, so $a>0$.  Since $X=-a$, equation \eqref{eq:norm-conclusion} yields
\begin{align*}
 X^3+X^2Y^2+Z^2+1
 &=-a^3+a^2Y^2+a^2(a-Y^2)\\
 &=0.
\end{align*}
This proves the lemma.
\end{proof}

Equivalently, \cref{lem:gaussian} is the norm identity obtained from
\[
 (P+iQ)^2(R+iS)=Z+i.
\]
We shall not need any theory of Gaussian integers beyond the explicitly expanded identity above.

\section{The polynomial identities}

For $t\in\Z$, define four integral linear polynomials
\begin{align}
 P(t)&=31487t+15586, \label{eq:P-poly}\\
 Q(t)&=-36963t-18295, \label{eq:Q-poly}\\
 R(t)&=-3400t-1682, \label{eq:R-poly}\\
 S(t)&=21114t+10451. \label{eq:S-poly}
\end{align}
The following proposition records the exact identities that make the construction work.

\begin{proposition}[Explicit polynomial identities]\label{prop:poly-identities}
For every integer $t$,
\begin{equation}\label{eq:identity-one}
 2P(t)Q(t)R(t)+\bigl(P(t)^2-Q(t)^2\bigr)S(t)=1,
\end{equation}
and
\begin{equation}\label{eq:identity-two}
 P(t)^2+Q(t)^2-R(t)^2-S(t)^2=F(t),
\end{equation}
where
\begin{equation}\label{eq:F-definition}
 F(t)=1900333542t^2+1881226506t+465577896.
\end{equation}
\end{proposition}

\begin{proof}
The verification is a direct calculation.  To make every cancellation auditable, we list the coefficient of each power of $t$.

For \eqref{eq:identity-one}, the two cubic polynomials on the left have the following coefficients:
\begin{center}
\small
\begin{tabular}{c r r r}
\toprule
power & $[t^j]\,2PQR$ & $[t^j]\,(P^2-Q^2)S$ & sum \\
\midrule
$j=3$ & $ 7914207070800$ & $-7914207070800$ & $0$ \\
$j=2$ & $11749892676484$ & $-11749892676484$ & $0$ \\
$j=1$ & $ 5814858098812$ & $-5814858098812$ & $0$ \\
$j=0$ & $  959230706680$ & $ -959230706679$ & $1$ \\
\bottomrule
\end{tabular}
\end{center}
Thus every nonconstant coefficient cancels and the constant coefficient is $1$, proving \eqref{eq:identity-one}.

For \eqref{eq:identity-two}, the corresponding quadratic coefficients are
\begin{center}
\small
\begin{tabular}{c r r r}
\toprule
power & $[t^j]\,(P^2+Q^2)$ & $[t^j]\,(R^2+S^2)$ & difference \\
\midrule
$j=2$ & $2357694538$ & $457360996$ & $1900333542$ \\
$j=1$ & $2333988934$ & $452762428$ & $1881226506$ \\
$j=0$ & $ 577630421$ & $112052525$ & $ 465577896$ \\
\bottomrule
\end{tabular}
\end{center}
This proves \eqref{eq:identity-two} and \eqref{eq:F-definition} coefficient by coefficient.
\end{proof}

Indeed, if $F(t)=y^2$ for some integers $t$ and $y$, then \eqref{eq:identity-one} and \eqref{eq:identity-two} are exactly the two hypotheses of \cref{lem:gaussian}.  The problem is therefore reduced to proving that $F(t)$ is a square for infinitely many integers $t$.  We give a fully self-contained Pell-type argument.

\section{A self-contained Pell propagation theorem}

For a fixed positive nonsquare integer $D$ and an element $\alpha=u+v\sqrt D$, define
\[
 \overline\alpha=u-v\sqrt D,
 \qquad
 \Norm(\alpha)=\alpha\overline\alpha=u^2-Dv^2.
\]
Direct multiplication gives $\overline{\alpha\beta}=\overline\alpha\,\overline\beta$ and hence
\begin{equation}\label{eq:norm-multiplicative}
 \Norm(\alpha\beta)=\Norm(\alpha)\Norm(\beta).
\end{equation}
We first prove the only existence fact about Pell equations that will be used.

\begin{lemma}[Existence of a nontrivial Pell unit]\label{lem:pell-existence}
Let $D>0$ be a nonsquare integer.  Then there exist integers $r>1$ and $s>0$ such that
\begin{equation}\label{eq:pell-basic}
 r^2-Ds^2=1.
\end{equation}
\end{lemma}

\begin{proof}
Since $D$ is a positive nonsquare integer, $\sqrt D$ is irrational.  For each integer $N\ge 2$, consider the $N+1$ fractional parts
\[
 \{0\sqrt D\},\{1\sqrt D\},\ldots,\{N\sqrt D\}
\]
in the interval $[0,1)$.  Partition $[0,1)$ into $N$ half-open intervals of length $1/N$.  By the pigeonhole principle, two of the fractional parts lie in the same interval.  Subtracting them, we obtain integers $p$ and $q$ such that
\begin{equation}\label{eq:dirichlet-approx}
 1\le q\le N,
 \qquad
 |q\sqrt D-p|<\frac1N.
\end{equation}
For all sufficiently large $N$, the integer $p$ is positive.

Infinitely many distinct pairs $(p,q)$ occur in this way.  Indeed, if only finitely many occurred, the finitely many positive numbers $|q\sqrt D-p|$ would have a positive minimum; this would contradict \eqref{eq:dirichlet-approx} for sufficiently large $N$.

For every such pair, put
\[
 k=p^2-Dq^2.
\]
The integer $k$ is nonzero, because $D$ is not a square and $q\ne0$.  Moreover, using \eqref{eq:dirichlet-approx}, $q\le N$, and
$p<q\sqrt D+1/N$, we obtain
\begin{align*}
 |k|
 &=|p-q\sqrt D|\,(p+q\sqrt D)\\
 &<\frac1N\left(2N\sqrt D+1\right)\\
 &<2\sqrt D+1.
\end{align*}
Thus $k$ belongs to a fixed finite set of nonzero integers.  One value of $k$ must therefore occur for infinitely many distinct positive pairs $(p,q)$.

Among those infinitely many pairs, choose two distinct pairs $(p_1,q_1)$ and $(p_2,q_2)$ that are componentwise congruent modulo $|k|$.  This is possible because there are only $|k|^2$ residue classes.  Define
\begin{equation}\label{eq:pell-rs-construction}
 r_0=\frac{p_1p_2-Dq_1q_2}{k},
 \qquad
 s_0=\frac{q_1p_2-p_1q_2}{k}.
\end{equation}
The congruences imply that both numerators in \eqref{eq:pell-rs-construction} are divisible by $k$, so $r_0,s_0\in\Z$.  Multiplying in $\Z[\sqrt D]$ gives
\begin{align*}
 r_0+s_0\sqrt D
 &=\frac{(p_1+q_1\sqrt D)(p_2-q_2\sqrt D)}{k}.
\end{align*}
Taking norms therefore yields
\[
 r_0^2-Ds_0^2
 =\frac{(p_1^2-Dq_1^2)(p_2^2-Dq_2^2)}{k^2}
 =1.
\]
We also have $s_0\ne0$.  Otherwise $q_1p_2=p_1q_2$, so the two positive pairs are positive rational multiples of one another.  Since both have the same nonzero norm $k$, the square of that multiplier is $1$, forcing the multiplier to be $1$ and the two pairs to be equal, a contradiction.

Finally, set $r=|r_0|$ and $s=|s_0|$.  Then $s>0$ and
$r^2=1+Ds^2>1$, so $r>1$.  Equation \eqref{eq:pell-basic} follows.
\end{proof}

We next arrange the congruences needed to return from a generalized Pell equation to an integral value of the original variable.

\begin{lemma}[Congruence-controlled Pell unit]\label{lem:congruence-unit}
Let $D>0$ be a nonsquare integer and let $M\ge1$.  There exist integers $\rho>1$ and $\sigma>0$ such that
\begin{align}
 \rho^2-D\sigma^2&=1, \label{eq:controlled-norm}\\
 \rho&\equiv1\pmod M, \label{eq:controlled-rho}\\
 \sigma&\equiv0\pmod M. \label{eq:controlled-sigma}
\end{align}
\end{lemma}

\begin{proof}
Choose $r>1$ and $s>0$ as in \cref{lem:pell-existence}.  If $M=1$, the congruence conditions are vacuous, and $(\rho,\sigma)=(r,s)$ is sufficient.  Hence assume $M\ge2$ and put
\[
 \varepsilon=r+s\sqrt D.
\]
Consider the finite ring
\[
 \mathcal R_M=(\Z/M\Z)[w]/(w^2-D).
\]
Because the defining polynomial is monic of degree $2$, every class has a unique representative $a+bw$ with $a,b\in\Z/M\Z$; in particular, $\mathcal R_M$ has exactly $M^2$ elements.  The residue class of $\varepsilon$ is a unit in $\mathcal R_M$, because its inverse is the residue class of $r-s\sqrt D$ and
\[
 (r+s\sqrt D)(r-s\sqrt D)=1.
\]
The unit group of the finite ring $\mathcal R_M$ is finite.  Hence the class of $\varepsilon$ has finite multiplicative order: for some $h\ge1$,
\[
 \varepsilon^h\equiv1\pmod M
\]
coefficientwise in the basis $1,\sqrt D$.  Write
\[
 \varepsilon^h=\rho+\sigma\sqrt D.
\]
Then \eqref{eq:controlled-rho} and \eqref{eq:controlled-sigma} hold.  Norms are multiplicative, so \eqref{eq:controlled-norm} holds as well.  Finally, because $r>1$ and $s>0$, every positive power of $r+s\sqrt D$ has positive integral coefficients; hence $\rho>1$ and $\sigma>0$.
\end{proof}

We can now prove the square-value theorem needed later.

\begin{proposition}[Pell propagation for quadratic square values]\label{prop:pell-propagation}
Let
\[
 f(T)=AT^2+BT+C\in\Z[T],
\]
where $A>0$ is not a square.  Assume that
\begin{equation}\label{eq:nonzero-discriminant}
 \Delta=B^2-4AC\ne0
\end{equation}
and that there is at least one pair $(T_0,Y_0)\in\Z^2$ satisfying
\begin{equation}\label{eq:initial-square}
 Y_0^2=f(T_0).
\end{equation}
Then $Y^2=f(T)$ has infinitely many integer solutions $(T,Y)$.  More precisely, the resulting integer values of $T$ are unbounded in absolute value.
\end{proposition}

\begin{proof}
Completing the square shows that \eqref{eq:initial-square} is equivalent to
\begin{equation}\label{eq:generalized-pell-initial}
 U_0^2-AV_0^2=\Delta,
\end{equation}
where
\begin{equation}\label{eq:UV-initial}
 U_0=2AT_0+B,
 \qquad
 V_0=2Y_0.
\end{equation}
Indeed,
\begin{align*}
 (2AT+B)^2-A(2Y)^2
 &=B^2+4A(AT^2+BT-Y^2),
\end{align*}
and this equals $B^2-4AC$ precisely when $Y^2=AT^2+BT+C$.

Apply \cref{lem:congruence-unit} with $D=A$ and $M=2A$.  We obtain $\rho>1$ and $\sigma>0$ such that
\begin{equation}\label{eq:prop-unit}
 \rho^2-A\sigma^2=1,
 \qquad
 \rho\equiv1\pmod{2A},
 \qquad
 \sigma\equiv0\pmod{2A}.
\end{equation}
For $n\ge0$, define $U_n,V_n\in\Z$ by
\begin{equation}\label{eq:pell-orbit}
 U_n+V_n\sqrt A=(U_0+V_0\sqrt A)(\rho+\sigma\sqrt A)^n.
\end{equation}
Equivalently, they satisfy the integral recurrence
\begin{align}
 U_{n+1}&=\rho U_n+A\sigma V_n, \label{eq:U-recurrence}\\
 V_{n+1}&=\sigma U_n+\rho V_n. \label{eq:V-recurrence}
\end{align}
Taking norms in \eqref{eq:pell-orbit} and using \eqref{eq:prop-unit} gives
\begin{equation}\label{eq:orbit-norm}
 U_n^2-AV_n^2=\Delta
 \qquad(n\ge0).
\end{equation}
The congruences in \eqref{eq:prop-unit}, together with \eqref{eq:U-recurrence} and \eqref{eq:V-recurrence}, imply inductively that
\begin{equation}\label{eq:orbit-congruences}
 U_n\equiv U_0\pmod{2A},
 \qquad
 V_n\equiv V_0\pmod{2A}.
\end{equation}
Since $U_0\equiv B\pmod{2A}$ and $V_0$ is even, the quantities
\begin{equation}\label{eq:TY-orbit}
 T_n=\frac{U_n-B}{2A},
 \qquad
 Y_n=\frac{V_n}{2}
\end{equation}
are integers.  Substituting \eqref{eq:TY-orbit} into \eqref{eq:orbit-norm}, and reversing the completed-square calculation, yields
\[
 Y_n^2=AT_n^2+BT_n+C=f(T_n).
\]
It remains to prove that infinitely many of these solutions are distinct.

Set
\[
 \alpha=U_0+V_0\sqrt A,
 \qquad
 \varepsilon=\rho+\sigma\sqrt A>1.
\]
Because $\Norm(\alpha)=\Delta\ne0$, we have $\alpha\ne0$.  The conjugate of $\varepsilon$ is $\varepsilon^{-1}$, since $\Norm(\varepsilon)=1$.  Hence \eqref{eq:pell-orbit} and its conjugate give
\begin{equation}\label{eq:Un-growth}
 U_n=\frac{\alpha\varepsilon^n+\overline\alpha\,\varepsilon^{-n}}{2}.
\end{equation}
The first term on the right of \eqref{eq:Un-growth} tends to infinity in absolute value, while the second tends to $0$.  Thus $|U_n|\to\infty$, and therefore $|T_n|\to\infty$ by \eqref{eq:TY-orbit}.  In particular, the pairs $(T_n,Y_n)$ include infinitely many distinct integer solutions.
\end{proof}

\begin{remark}\label{rem:why-discriminant}
The hypothesis $\Delta\ne0$ in \cref{prop:pell-propagation} is used only to ensure that the initial generalized-Pell element $U_0+V_0\sqrt A$ is nonzero, which makes the orbit unbounded.  In the application below, the discriminant is explicitly negative and nonzero.
\end{remark}

\section{Infinitely many square values of the auxiliary quadratic}

We apply \cref{prop:pell-propagation} to the polynomial $F$ in \eqref{eq:F-definition}.  Set
\begin{equation}\label{eq:ABC-values}
 A=1900333542,
 \qquad
 B=1881226506,
 \qquad
 C=465577896.
\end{equation}
First, $A$ is not a square, because
\begin{equation}\label{eq:A-nonsquare}
 43592^2=1900262464
 <1900333542
 <1900349649=43593^2.
\end{equation}
Second, the discriminant is
\begin{equation}\label{eq:F-discriminant}
 B^2-4AC=-1853382492\ne0.
\end{equation}
Third, $F$ has the following integral square value:
\begin{align}
 F(-659)
 &=1900333542\cdot659^2-1881226506\cdot659+465577896 \notag\\
 &=825278750953302-1239728267454+465577896 \notag\\
 &=824039488263744 \notag\\
 &=28706088^2. \label{eq:initial-F-square}
\end{align}
All hypotheses of \cref{prop:pell-propagation} are therefore satisfied.

\begin{corollary}\label{cor:F-infinitely-many-squares}
There exist infinitely many pairs $(t,y)\in\Z^2$ satisfying
\begin{equation}\label{eq:F-square}
 y^2=F(t).
\end{equation}
Moreover, the corresponding integers $t$ are unbounded in absolute value.
\end{corollary}

\begin{proof}
Apply \cref{prop:pell-propagation} using \eqref{eq:ABC-values}, \eqref{eq:A-nonsquare}, \eqref{eq:F-discriminant}, and \eqref{eq:initial-F-square}.
\end{proof}

For completeness, the Pell orbit in this application begins with
\begin{align*}
 U_0&=2A(-659)+B=-2502758381850,\\
 V_0&=2\cdot28706088=57412176,
\end{align*}
and it satisfies
\[
 U_0^2-AV_0^2=-1853382492.
\]
Any norm-one Pell unit $\rho+\sigma\sqrt A$ satisfying
\[
 \rho\equiv1\pmod{2A},
 \qquad
 \sigma\equiv0\pmod{2A}
\]
then gives the explicit recurrence \eqref{eq:U-recurrence}, \eqref{eq:V-recurrence}, and \eqref{eq:TY-orbit}.  The existence of such a unit was proved in \cref{lem:congruence-unit}; no unproved appeal to the classical Pell theorem remains.

\section{Proof of the main theorem}

\begin{proof}[Proof of \cref{thm:main}]
Let $(t,y)$ be any of the infinitely many integer pairs supplied by \cref{cor:F-infinitely-many-squares}.  Define $P=P(t)$, $Q=Q(t)$, $R=R(t)$, and $S=S(t)$ by equations \eqref{eq:P-poly}--\eqref{eq:S-poly}.  By \cref{prop:poly-identities},
\[
 2PQR+(P^2-Q^2)S=1
 \quad\text{and}\quad
 P^2+Q^2-R^2-S^2=F(t)=y^2.
\]
Therefore the hypotheses of \cref{lem:gaussian} hold.  Define
\begin{align}
 x(t)&=-\bigl(P(t)^2+Q(t)^2\bigr), \label{eq:x-family}\\
 z(t)&=\bigl(P(t)^2-Q(t)^2\bigr)R(t)-2P(t)Q(t)S(t). \label{eq:z-family}
\end{align}
For reference, the latter polynomial expands to
\begin{equation}\label{eq:z-expanded}
 z(t)=50421655389668t^3+74872031013786t^2
      +37059612550518t+6114499038718.
\end{equation}
Then
\[
 x(t)^3+x(t)^2y^2+z(t)^2+1=0.
\]
In fact, the same calculation before specializing $F(t)$ to the square $y^2$ gives the polynomial identity
\begin{equation}\label{eq:master-polynomial-identity}
 x(t)^3+x(t)^2F(t)+z(t)^2+1=0
 \qquad\text{in }\Z[t].
\end{equation}

It remains only to verify that infinitely many distinct triples arise.  From the coefficient calculation in \cref{prop:poly-identities},
\begin{equation}\label{eq:a-polynomial}
 P(t)^2+Q(t)^2
 =2357694538t^2+2333988934t+577630421.
\end{equation}
This quadratic has positive leading coefficient, so it tends to $+\infty$ as $|t|\to\infty$.  By \cref{cor:F-infinitely-many-squares}, the admissible values of $t$ are unbounded in absolute value.  Consequently, \eqref{eq:a-polynomial} assumes infinitely many values on those admissible integers, and hence the integers $x(t)$ in \eqref{eq:x-family} assume infinitely many distinct values.  The resulting solution triples are therefore pairwise distinct along an infinite subsequence.
\end{proof}



\section{Concrete solutions and independent checks}

Taking the initial Pell point $t=-659$ gives
\begin{align*}
 P(-659)&=-20734347,
 &Q(-659)&=24340322,\\
 R(-659)&=2238918,
 &S(-659)&=-13903675,
\end{align*}
and \eqref{eq:initial-F-square} gives $y=\pm28706088$.  The definitions \eqref{eq:x-family} and \eqref{eq:z-family} yield
\begin{align*}
 x&=-1022364420580093,\\
 z&=-14397741918770263093350.
\end{align*}
Thus, for example,
\begin{equation}\label{eq:large-example}
 \boxed{
 (x,y,z)=
 (-1022364420580093,\ 28706088,\ -14397741918770263093350)
 }
\end{equation}
is an integer solution.  Its validity follows from the proved identities, but it may also be checked directly by exact integer arithmetic.

There is also a much smaller solution not needed for the infinite-family proof:
\begin{equation}\label{eq:small-example}
 (x,y,z)=(-493,12,9210).
\end{equation}
Indeed,
\[
 493-12^2=349
 \quad\text{and}\quad
 9210^2+1=84824101=493^2\cdot349.
\]
Hence
\[
 (-493)^3+(-493)^2\cdot12^2+9210^2+1=0.
\]
By sign symmetry, both signs of $12$ and both signs of $9210$ give solutions.

\section{Conclusion}

The proof supplies an unconditional infinite construction.  The essential points are:
\begin{enumerate}
 \item the two-square identity in \cref{lem:gaussian}, which converts the conditions \eqref{eq:bezout-condition} and \eqref{eq:square-difference} into a solution of the original equation;
 \item the exact linear-polynomial identities in \cref{prop:poly-identities}; and
 \item the self-contained Pell propagation theorem, which forces the quadratic $F(t)$ to be a square for infinitely many unbounded integer values of $t$.
\end{enumerate}
Therefore \eqref{eq:original} has infinitely many integer solutions, as claimed.

\end{document}
